% Transcription — fonds Alexandre Grothendieck, shelfmark 55, batch 1 (pages 1-17).
% Established from the facsimile archives/batches/55/batch-01.pdf (a mirror of
% https://grothendieck.umontpellier.fr/55.pdf), per the skill
% transcribe-grothendieck. The folder is seventeen pages and this is its only
% batch. Sixteen of the seventeen carry his hand and are transcribed; page 15
% carries nothing and is skipped, the gap in the page numbers being the record.
%
% The folder holds two things. Pages 1 to 14 alternate a leaf of manuscript
% notes, in blue ink, with a leaf of typescript that he annotated: pages 1, 3,
% 5, 7, 9, 11 and 13 are the notes, pages 2, 4, 6, 10, 12 and 14 are pages
% 111, 110, 109, 101, 102 and 104 of a typed chapter on dimension and
% codimension in topological spaces and in locally noetherian preschemes, and
% page 8 is page « 8 » of another typescript, on associated prime cycles,
% typed on a machine without accents and referring to « IV 2.8.18 ». The
% typescript leaves are kept because they carry his hand — insertions,
% substitutions and marginal notes, some of them substantial (« propositions
% laborieuses inutiles », « donner le cas d'égalité par platitude », the
% systematic replacement of « points fermés de même codimension » by
% « équidimensionnel » in Prop. 5 of page 14) — and, following folder 48, the
% statements they carry and his interventions are given while the typed proofs
% are summarised in notes rather than recomposed. The typed pages bear no
% relation to the manuscript notes they alternate with. Pages 16 and 17 are a
% separate fragment in black ink and a steady hand, with no title and nothing
% tying it to the rest; its section heading is the edition's.
%
% The inventory's title « Espaces analytiques » is not written anywhere in the
% folder. What the manuscript notes are about is narrower: page 1 opens on
% « X espace de Stein compact, A = Γ(X, O_X) » and proves that
% M ↦ M ⊗_{A_X} O_X is an equivalence between A-modules of finite presentation
% and coherent O_X-Modules — full faithfulness from a presentation
% A^p → A^q → M → 0 and the Stein vanishing, essential surjectivity from
% global generation (page 3, with its commutative diagram). Page 5 abstracts
% this into Lemme A for a ringed space under hypotheses (i) quasi-compact,
% (ii) O_X coherent with noetherian stalks, (iii) H^1(X, F) = 0 for coherent
% F, with a variant written up the left margin, three corollaries (exactness
% and so flatness of O_{X,x} over A; A is « cohérent »; a third on maximal
% ideals, largely illegible) and a Remarque replacing (ii)-(iii) by
% (α) exactness of Γ on presentations and (β) global generation. Page 7 draws
% the consequences for algebras of finite presentation and asks, as « une
% question de pure algèbre commutative », whether the kernel of
% O_X → B is coherent, noting « ok. dans le cas classique X polycylindre
% analytique ». Page 11 (Corollaire, conditions (ii') and (iii')) extends
% Lemme A to O_X[T_1, ..., T_n] and to any algebra of finite presentation by
% writing a presentation as a limit of presentations by coherent O_X-Modules
% of bounded degree, and page 9 reads as its continuation — the archivists'
% order of pages 9 and 11 is probably not the order of writing, which is
% flagged where it occurs and not repaired — ending on the coherence of the
% J_m = J ∩ B_m and the marginal « ? histoire d'anneaux cohérents ? ».
% Page 13 is Lemme B: under (i), (ii), (iv) noetherian stalks and (v) an
% identity principle on closed analytic subspaces with integral local ring,
% increasing sequences of coherent subsheaves are stationary, whence the
% corollary that A = Γ(X, O_X) is noetherian, with the struck « espace de
% Stein » replaced by « polycylindre analytique (compact) » and the circled
% marginal « premier espoir ». Pages 16 and 17 show that there is a complex
% torus of dimension 2, fibred over a complex torus of dimension 1 with a
% one-dimensional subtorus as fibre, that is not an abelian variety — T is
% C^2 modulo Γ generated by e_1, ie_1, e_2, e_1 + ike_2 with k irrational —
% by an explicit computation with the components of x and cx in Γ.
%
% The manuscript hand is drafting throughout: the statements, the displayed
% formulas and the underlined words carry, the connective prose often does
% not, and pages 5, 9 and 13 are the worst of it, each with blocks boxed or
% cancelled over a crowded sheet. Following folder 29 batches 8 and 9 and the
% passes since, what the leaves will not support is left \ill{} rather than
% invented. The marginal notes on the typescript leaves are written in biais,
% up the left edge, and have been read turned; the long ones on pages 2 and 4
% are given only as far as the strokes carry.
%
% Notation fixed once here rather than page by page. He underlines the O of
% the structure sheaf, and it is set \underline{O}_X; in the typescripts the
% typed O is set \mathcal{O}. « A_X » in M ⊗_{A_X} O_X is the constant sheaf of
% rings, and M~ is M ⊗_{A_X} O_X. « prés. finie » and « pl. fidèle » are his
% abbreviations for « présentation finie » and « pleinement fidèle »;
% « loc noeth », « f et s », « pt gén », « comp irr » are the typescript's.
% He writes a script B for an O_X-algebra and Γ B for its sections.
%
% Readings flagged and not settled: « sur X » in the statement 1°) of page 1,
% where A is what the sense wants; the adjective « globale » qualifying
% « présentation finie » twice on page 7; « pur » in Corollaire 3 of page 5
% and in the cancelled N.B. of page 9; the word « \Gamma-propre » on page 16;
% « troisièmes » on page 17; and the two last members of the equality in the
% marginal corollary of page 14.
%
% Two variances of substance are flagged where they occur and not repaired.
% Page 16 writes « V soit Γ-propre, i.e. V engendré par Γ ∩ V » where the
% argument that follows is about the supplement V'. And the marginal inequality
% at the foot of page 4, dim O_x ≤ dim O_y + dim O_x, is read as written though
% its last term does not carry the argument the sense wants.
%
% Pass: Opus 5 (claude-opus-5), 2026-09-13 — first pass, unchecked against the pages by a human.
\documentclass[11pt,a4paper]{article}
% Le préambule commun à toutes les transcriptions du fonds Grothendieck.
%
% Il tient en une idée : **l'incertitude se compose, elle ne se résout pas.**
% Une transcription qui lisse un mot douteux a détruit la seule chose pour
% laquelle on la faisait — un lecteur doit pouvoir savoir, à chaque endroit,
% ce qui était sur le feuillet et ce qui a été ajouté. Les six macros qui
% suivent sont donc l'appareil critique, et non de la décoration.
%
% Les mêmes commandes sont reconnues par `scripts/render.mjs`, qui produit la
% vue de lecture du site. Une macro ajoutée ici doit l'être aussi là-bas,
% faute de quoi le rendu échouera bruyamment — ce qui est voulu : un rendu
% silencieusement faux est pire qu'un rendu absent.

\ProvidesPackage{grothendieck}[2026/08/08 Transcriptions du fonds Grothendieck]

\usepackage{amsmath,amssymb,amsthm}
\usepackage{mathrsfs}
\usepackage{stmaryrd}
% Les diagrammes commutatifs sont partout dans ce fonds : c'est la langue
% même de la géométrie algébrique de Grothendieck, et une transcription qui
% les rendrait en prose ne transcrirait rien.
\usepackage{tikz-cd}
% Français par défaut — c'est la langue des feuillets — mais l'anglais est
% chargé aussi : la traduction et le résumé sont des documents anglais, et la
% typographie française (espaces avant les deux-points) y serait une faute.
% Ces documents basculent par \selectlanguage{english} en tête de corps.
\usepackage[english,french]{babel}
\usepackage{fontspec}
\usepackage{geometry}
\geometry{a4paper,margin=25mm,marginparwidth=28mm}
\usepackage{marginnote}
\usepackage{xcolor}
% ulem, not soul. `soul`'s \st and \ul are built on a character-by-character
% scan that XeLaTeX's font machinery breaks: the document fails with an
% undefined control sequence rather than degrading. ulem does the same two jobs
% and is Unicode-engine safe. `normalem` stops it hijacking \emph, which the
% transcriptions use for Grothendieck's own emphasis.
\usepackage[normalem]{ulem}
\usepackage{hyperref}
\hypersetup{colorlinks=true,linkcolor=black,urlcolor=black,pdfborder={0 0 0}}

% --- Métadonnées du lot -----------------------------------------------------
% Reprises telles quelles par le script de rendu, qui les affiche en tête de
% la vue de lecture. Elles fixent ce que le fichier prétend couvrir : sans
% elles, une transcription flotte sans point d'attache dans un fonds de seize
% mille feuillets.

\newcommand{\folder}[1]{\gdef\grothFolder{#1}}
\newcommand{\batch}[1]{\gdef\grothBatch{#1}}
\newcommand{\leaves}[2]{\gdef\grothFirst{#1}\gdef\grothLast{#2}}
\newcommand{\foldertitle}[1]{\gdef\grothFoldertitle{#1}}
\gdef\grothFolder{?}\gdef\grothBatch{?}\gdef\grothFirst{?}\gdef\grothLast{?}\gdef\grothFoldertitle{}

% --- Le résumé --------------------------------------------------------------
%
% Il ouvre la lecture modernisée, sous le titre « Résumé », et il est composé
% en petit. Ce n'est pas un ornement : c'est le seul passage du document qui
% ne soit pas des mathématiques, et un lecteur doit pouvoir le reconnaître
% avant de l'avoir lu — puis le sauter s'il connaît déjà le sujet, ou s'y
% arrêter s'il ne le connaît pas. Le filet qui le suit marque la reprise du
% corps.

\newenvironment{resume}
  {\small\setlength{\parskip}{0.45em}}
  {\par\medskip\hrule\medskip}

% --- Mots-clés ---------------------------------------------------------------
%
% En anglais, en clôture du résumé de la lecture modernisée : le vocabulaire
% moderne sous lequel chercher ce que les feuillets construisent. Le manifeste
% du site (`npm run manifest`) les extrait pour étiqueter la cote entière —
% cette ligne est la source unique des tags, il n'y a pas de fichier à côté.
\newcommand{\keywords}[1]{\par\smallskip\noindent
  {\footnotesize\textsc{Keywords} --- \textit{#1}}\par}

% --- L'appareil critique ----------------------------------------------------

% Le numéro de feuillet, en marge. C'est l'ancre qui permet de régler un
% désaccord sur une formule en regardant *un* feuillet plutôt qu'une chemise
% de six cents. Le site s'en sert aussi pour tourner les pages du fac-similé
% quand on fait défiler la transcription.
\newcommand{\leaf}[1]{%
  \marginnote{\footnotesize\color{gray!70}\textsf{#1}}%
  \label{leaf:#1}\ignorespaces}

% Illisible : on ne devine pas. Un mot inventé qui se lit comme les autres est
% le pire résultat possible.
\newcommand{\ill}{\textcolor{red!60!black}{[\,\ldots\,]}}

% Lecture douteuse : le mot est donné, et signalé comme incertain.
\newcommand{\uncertain}[1]{\uline{#1}}

% Ajout de l'éditeur — une lettre manquante, un mot sous-entendu.
\newcommand{\add}[1]{\textcolor{blue!50!black}{[#1]}}

% Biffé par Grothendieck lui-même. Ce qu'il a rayé fait partie du document :
% une rature dit souvent où la pensée a changé de direction.
\newcommand{\struck}[1]{\sout{#1}}

% Note du transcripteur, sur l'état matériel du feuillet.
\newcommand{\note}[1]{\par\smallskip{\small\color{gray!60!black}\textit{[#1]}}\par\smallskip}

% Note marginale de Grothendieck, distincte du corps du texte.
\newcommand{\marginal}[1]{\par\smallskip\noindent\hspace{1em}%
  {\small\color{gray!60!black}#1}\par\smallskip}

% --- Titre ------------------------------------------------------------------

\AtBeginDocument{%
  \begin{center}
    {\large\bfseries Fonds Alexandre Grothendieck --- cote n\textsuperscript{o}~\grothFolder}\\[2pt]
    {\itshape\grothFoldertitle}\\[4pt]
    {\small feuillets \grothFirst--\grothLast\ (lot \grothBatch)}\\[2pt]
    {\footnotesize\color{gray!60!black}Transcription établie d'après le fac-similé numérisé
     par l'Université de Montpellier.\\
     Les lectures incertaines sont soulignées, les illisibles notées [\,\ldots\,],
     les ajouts entre crochets.}
  \end{center}
  \vspace{1em}}

\endinput


\folder{55}
\batch{1}
\pages{1}{17}
\dating{[à partir de 1961]}
\watermark{Édition de démonstration}
\foldertitle{Espaces analytiques : notes manuscrites (s.d.)}

\begin{document}

\section*{Espaces de Stein et Modules cohérents}

\note{titre de l'édition : les notes ne portent pas de titre, et celui de
l'inventaire, « Espaces analytiques », n'est écrit nulle part dans le dossier}

\page{1}

$X$ espace de Stein compact

$A = \Gamma(X, \underline{O}_{X})$

\struck{$B \in$ Alg. de présentation finie}

1°) Le foncteur $M \mapsto M \otimes_{A_{X}} \underline{O}_{X}$ de la catégorie
des modules de \struck{type} présentation finie \uncertain{sur} $X$\note{la
lettre est bien un $X$, là où le sens attendrait $A$} dans la catégorie des
$\underline{O}_{X}$-Modules \emph{coh.} est une équivalence de catégories.

\struck{En effet, le foncteur est exact (?) pour voir s'il est pl. fid. il suffit
de prouver que $\widetilde{M} = 0 \Rightarrow M = 0$. Or}\note{trois lignes
encadrées et biffées d'un trait}

Soit $M$ de présentation finie,
\[
A^{p} \to A^{q} \to M \to 0
\]
d'où une \ill{}
\[
\underline{O}^{p} \to \underline{O}^{q} \to \widetilde{M} \to 0
\]
d'où prenant les $\Gamma$, et notant l'exactitude (car $X$ Stein)
\[
\struck{\Gamma(M)} \quad M \xrightarrow{\sim} \Gamma(\widetilde{M})
\]
Donc \struck{\ill{}} \uncertain{tout}
\[
\struck{\widetilde{M} = 0 \Rightarrow \Gamma(\widetilde{M}) = 0 \text{ i.e. } M = 0}
\]

\page{2}

\note{feuillet de tapuscrit paginé « -- 111 -- », repris à la
main. Il appartient à un exposé sur la dimension et la codimension
dans les préschémas localement noethériens, sans rapport avec les espaces de
Stein de la page 1 ; il est gardé parce qu'il porte sa main. Comme pour les
tapuscrits du dossier 48, on donne les énoncés que porte le feuillet et ses
interventions, et le corps du tapuscrit n'est pas re-composé}

\marginal{propositions laborieuses inutiles}\note{en regard des quatre
premières lignes, que ce trait vertical embrasse : la fin d'une démonstration où
l'on se ramène à $Y = \mathrm{Spec}(B)$, $B$ anneau local \emph{noethérien},
$X = \mathrm{Spec}(A)$, $A$ algèbre \emph{noethérienne} sur $B$ --- les deux
adjectifs sont portés en interligne de sa main au-dessus d'un mot tapé
biffé --- et où l'on montre que $\mathrm{rg}(\mathfrak{p}) \leq
\mathrm{rg}(\mathfrak{m})$ pour $\mathfrak{p}$ premier minimal associé à
$A\mathfrak{m}$}

\textbf{Corollaire 1.} Soient \add{$X$ et} $Y$ \struck{\ill{}} \add{deux}
préschémas localement noethériens, $f : X \to Y$ un morphisme \struck{\ill{}},
$Y'$ une partie fermée irréductible de $Y$, $X'$ une composante irréductible de
$f^{-1}(Y')$, dont l'image dans $Y$ soit dense dans $Y'$. Alors
$\mathrm{codim}_{X}(X') \leq \mathrm{codim}_{Y}(Y')$.\note{« $X$ et » et
« deux » sont insérés à la main au-dessus de la ligne, un mot tapé étant biffé
entre eux}

\textbf{Corollaire 2.} Soient $X$ un préschéma \add{localement noethérien}
\struck{\ill{}}, $f$ une section de $\mathcal{O}_{X}$ au-dessus de $X$, qui
n'est identiquement nulle dans aucune composante irréductible de $X$. Soit
$Z_{f}$ l'ensemble des $x \in X$ tels que $f(x) = 0$. Alors toute composante
irréductible de $Z_{f}$ est de codimension $1$.\note{« localement noethérien »
en interligne, de sa main, au-dessus d'un mot tapé biffé. La démonstration
tapée se ramène au cas affine noethérien et invoque le « Hauptidealsatz » de
Krull (cf. Bourbaki, Alg. comm., la référence restant en blanc)}

\marginal{affirmation\add{s} inutile\add{s}, si on remarque que \ill{} $x \in X$
\ill{} \ill{} \ill{} \ill{} \ill{} $2$ \ill{} \ill{} \ill{}}\note{note écrite
de biais, à quelque soixante degrés, dans la marge gauche en regard de la
démonstration du corollaire 2 ; seul son début se laisse lire}

\page{3}

Soit $N$ un autre module de prés. finie, on a le diagramme commutatif à lignes
exactes

\begin{tikzcd}
0 \arrow[r] & \mathrm{Hom}(M, N) \arrow[r] \arrow[d, "\alpha"] & N^{q} \arrow[r] \arrow[d, "\beta"] & N^{p} \arrow[d, "\gamma"] \\
0 \arrow[r] & \mathrm{Hom}(\widetilde{M}, \widetilde{N}) \arrow[r] & \Gamma\widetilde{N}^{q} \arrow[r] & \Gamma\widetilde{N}^{p}
\end{tikzcd}

et comme $\beta$ et $\gamma$ sont bijectives, \uncertain{on a} $\alpha$
\ill{}, donc le foncteur $M \mapsto \widetilde{M}$ est \emph{pleinement
fidèle}. Prouvons qu'il est essentiellement surjectif : en effet, si $F$
\ill{} coh. sur $X$, on sait qu'il est engendré par ses sections, i.e.
$\exists\ \underline{O}_{X}^{q} \to F \to 0$ \struck{\ill{}}
\uncertain{exact}.\note{le dernier mot est porté en interligne au-dessus d'un
mot biffé}

Soit $R$ le noyau, alors il existe de même $\underline{O}_{X}^{p} \to R \to 0$,
d'où une suite exacte
\[
\underline{O}_{X}^{p} \xrightarrow{u} \underline{O}_{X}^{q} \to F \to 0
\]
\struck{et par suite \ill{} \ill{} \ill{} \ill{}} or $u$ définit (grâce aux
sections)
\[
A^{p} \xrightarrow{u} A^{q}
\]
\ill{} \uncertain{prenant} $M = \mathrm{Coker}\, u$ \ill{} \ill{} \ill{} par
\[
F \simeq \widetilde{M}.
\]

\page{4}

\note{feuillet du même tapuscrit, paginé « -- 110 -- » : il précède, dans la pagination du tapuscrit, le feuillet donné en page 2.
Il porte la fin de la démonstration de $\dim(X \times_{k} Y) = \dim(X) +
\dim(Y)$ pour des préschémas de type fini sur un corps $k$ --- réduction au cas
affine intègre, lemme de normalisation, théorème de Cohen-Seidenberg --- puis
l'énoncé de la proposition 8}

\marginal{abusif}\note{en regard de la première ligne : « $\dim(X \times_{k} Y)
= \dim(X) + \dim(Y)$ ; \emph{en particulier, pour tout corps $K$, extension de
$k$, $\dim(X \otimes_{k} K) = \dim(X)$} », les mots « en particulier » étant
soulignés à la main d'un trait ondulé}

\marginal{\ill{} \ill{} $K(z)$ ($z \in X \times_{k} Y$) \ill{} \ill{} \ill{}
\ill{} calc. \ill{} \ill{} de l'ext. $K(x)$ et $K(y)$ ($x \in X$, $y \in Y$)
\ill{} \ill{} \ill{} \ill{}}\note{note écrite de biais le long de la marge
gauche, en regard du lemme de normalisation ; « sur $k$ » est d'autre part
inséré à la main après « le degré de transcendance » du corps des fractions de
$C/\mathfrak{p}$}

\textbf{Proposition 8.} Soient \add{$X$,} $Y$ \struck{un} \add{deux} préschémas
localement noethériens, $f : X \to Y$ un morphisme \struck{de type fini}.
Soient $y \in Y$, et soit $x$ le point générique d'une composante irréductible
de $f^{-1}(y)$, considéré comme préschéma sur $\kappa(y)$ (§ 2, n° 7, prop. 17).
Alors $\mathrm{codim}_{X}\overline{\{x\}} \leq
\mathrm{codim}_{Y}\overline{\{y\}}$ (ou, ce qui revient au même, $\dim
\mathcal{O}_{x} \leq \dim \mathcal{O}_{y}$).\note{« $X$, » et « deux » en
interligne, « un » et « de type fini » biffés, de sa main ; « La question est
locale en $x$ et en $y$ », qui suit, est la première ligne de la démonstration}

\marginal{Simplifier dans la preuve \ill{} $\dim \mathcal{O}_{x} \leq \dim
\mathcal{O}_{y} + \dim \mathcal{O}_{x}$ \ill{}}\note{dans la marge inférieure
gauche, avec une flèche qui renvoie à trois lignes serrées écrites sous le
dernier alinéa tapé et barrées de longs traits obliques ; on y distingue
$f^{-1}(y)$ et rien de suivi. La dernière inégalité est lue telle quelle, son
dernier terme ne portant pas l'argument que le sens demande}

\marginal{donner le cas d'égalité par \emph{platitude}}

\page{5}

\note{le lemme A qui ouvre la page est entouré d'un crochet ouvrant et d'un
crochet fermant ; il reprend en l'abstrayant l'énoncé 1°) de la page 1}

\textbf{Lemme A}\quad $X$ espace annelé, \struck{avec $\underline{O}_{X}$
cohérent}\note{au-dessus du passage biffé, deux mots en interligne que les
traits ne laissent pas lire}
\begin{itemize}
\item[(i)] $X$ quasi-compact
\item[] \struck{\ill{} $\underline{O}_{X}$ cohérent}
\item[(ii)] $\underline{O}_{X}$ cohérent, les $\underline{O}_{X,x}$
noethériens.
\item[(iii)] Pour \struck{\ill{} \ill{} \ill{}} tout Module \struck{cohérent}
\uncertain{cohérent} $F$ sur $X$, on a $H^{1}(X, F) = 0$ \struck{dans les cas
suivants :} \struck{\ill{} \ill{} loc. de type fini sur \ill{}}\note{« (iii) »
est réécrit dans un cercle en regard d'un premier numéro surchargé ; « tout »
et le second « cohérent » sont portés en interligne}
\end{itemize}
Alors, posant $A = \Gamma(X, \underline{O}_{X})$, le foncteur
\[
M \mapsto \widetilde{M} = M \otimes_{A_{X}} \underline{O}_{X}
\]
de la catégorie des $A$-modules de présentation finie dans la catégorie des
$\underline{O}_{X}$-Modules de prés. finie (= cohérents) est une équivalence
de catégories.

\marginal{\emph{Variante}. Si tout pt de $X$ a un système fond. de voisinages
\struck{\ill{}} quasi-compacts, si $X$ étant quasi-compact, et si pour tout
\ill{} Module $F$ \ill{} \ill{} image d'un $\underline{O}^{p} \to
\underline{O}^{q}$, $H^{1}(X, F) = 0$, alors le foncteur $M \mapsto M
\otimes_{A_{X}} \underline{O}_{X}$ des modules sur $A$ dans les Modules
\struck{\ill{}} \ill{} \ill{} \ill{} $\underline{O}^{I} \to
\underline{O}^{J}$, \ill{} \struck{\ill{}} \ill{} équivalence \ldots}\note{note
écrite verticalement dans la marge gauche, de bas en haut, sur toute la hauteur
du lemme A, et qu'une série de traits verticaux sépare du corps ; « de
présentation finie », en interligne au-dessus de sa dernière ligne, est
biffé. Sous elle, quelques traits et signes isolés qui ne se lisent pas}

\textbf{Corollaire 1}\quad Le foncteur $M \mapsto \widetilde{M}$ est
\emph{exact}. En particulier, pour tout $x \in X$, $\underline{O}_{X,x}$ est
\emph{plat} sur $A$.

\textbf{Corollaire 2}\quad L'anneau $A$ est « cohérent » i.e. pour tout hom
$A^{n} \to A$, le noyau est de type fini.

\textbf{Corollaire 3}\quad Si $X$ est « \uncertain{pur} », \struck{$\underline{O}_{X}$}
\ill{} on a \ill{} \ill{} \ill{} idéaux maximaux de \struck{type} \uncertain{fini}
de $A$, et \ill{} \ill{} maximal de $\underline{O}_{X,x}$.\note{la ligne
suivante, écrite serrée sous l'énoncé et en partie biffée, porte « $(x,
\mathfrak{m})$ », « $\mapsto x$ » et « idéal » et ne se laisse pas lire
davantage}

\emph{Remarque}\quad on peut \uncertain{remplacer} (ii) (iii) par
\begin{itemize}
\item[($\alpha$)] Pour toute suite exacte $\underline{O}_{X}^{p} \to
\underline{O}_{X}^{q} \to F \to 0$, \ill{} exactitude pour la suite des
$\Gamma$
\item[($\beta$)] Tout $F$ de présentation finie est engendré par ses
sections \ldots
\end{itemize}

\page{6}

\note{feuillet du même tapuscrit, paginé « -- 109 -- ».
Il porte la fin d'un contre-exemple, les corollaires 2 et 3 d'une proposition
sur la dimension des préschémas algébriques sur un corps, un second
contre-exemple et le début de l'énoncé du corollaire 4 ; on donne les deux
énoncés et ce qu'il y ajoute}

\textbf{Corollaire 2.} Soit $X$ un préschéma algébrique sur un corps $k$. Pour
tout point $x \in X$, on a $\dim_{x}(X) = \sup_{i}(\dim(X_{i}))$, où $X_{i}$
parcourt l'ensemble des composantes irréductibles de $X$ contenant $x$. Si
$\{x\}$ est fermé on a $\dim_{x} X = \mathrm{codim}_{X} x = \dim(\mathcal{O}_{x})$,
\add{et en tout cas, on a} $\dim_{x} X = \dim \underline{O}_{X,x} + \dim
k(x)/k$.\note{la fin, après la virgule, est ajoutée à la main à la suite de la
ligne tapée. Dans la démonstration, il entoure d'un trait « pour tout ouvert
$U$ contenant $x$ » et insère « si $\{x\}$ est fermé » après « d'après le
cor. 1, $\dim_{x}(X_{i}) = \dim(X_{i})$ » ; une petite ligne à l'encre surmonte
en outre le « $\sup_{i}$ » de l'énoncé}

\textbf{Corollaire 3.} Soient $X$, $Y$ deux préschémas algébriques sur un corps
$k$. Pour tout $k$-morphisme $f : X \to Y$, on a $\dim(\overline{f(X)}) \leq
\dim(X)$.

\marginal{\uncertain{a-t-on} \ill{} \ill{} \struck{\ill{}} \struck{dans le}
\struck{cas} \struck{\ill{}}}\note{en regard de l'énoncé du corollaire 3 ; les
derniers mots, en crayon plus pâle, sont barrés un à un et un double trait
vertical les borde. Dans la démonstration, « alors » est inséré à la main
au-dessus de « telle que »}

\marginal{Corollaire \quad $\dim_{x} X = \dim \underline{O}_{x} + \dim
k(x)/k$ \ill{} \ill{} \ill{}}\note{dans la marge inférieure gauche, en regard du
second contre-exemple ($B$ anneau de valuation discrète, $X = \mathrm{Spec}(K)
\to Y = \mathrm{Spec}(B)$ de type fini, $\dim(Y) = 1$ et $\dim(X) = 0$) ; la
formule est écrite en deux morceaux reliés par un trait sinueux, et les deux
derniers mots, sous un « = » isolé, ne se lisent pas}

\page{7}

Du lemme on tire \ill{}, sans utiliser \uncertain{l'hyp.} ($\alpha$), \ill{}
\ill{} \ill{}, les conclusions :

$B \mapsto B \otimes_{A_{X}} \underline{O}_{X}$ est une équivalence de la
catégorie des $A$-Alg. de présentation finie \ill{} la catégorie des
$\underline{O}_{X}$-Alg. de présentation \emph{\uncertain{globale}} finie.

Si $B$ \ill{} \ill{} $\mathcal{B}$ la \ill{}, alors
\[
M \mapsto M \otimes_{B_{X}} \mathcal{B}
\]
est une \ill{} \ill{} catégorie des $B$-Modules de présentation finie, et la
catégorie des $\mathcal{B}$-Modules de présentation \emph{\uncertain{globale}}
finie.

\emph{Question}\quad Soit $\mathcal{B}$ \add{Alg.} de prés. finie sur $X$, le
noyau $J$ de $\underline{O}_{X}$ \ill{} $\to \mathcal{B}$ est-il cohérent ??
i.e. c'est une question de pure algèbre commutative, peut s'énoncer : $A$
\ill{} un anneau cohérent, \struck{$A$} $J$ un idéal de type fini dans $B =
A[T]$, \ill{} \ill{} $I$ \ill{} $A$ \ill{} de type fini ?? [Plus général, est-il
vrai que \ill{} $I \subset B$ sous-$A$-modules de type fini ??]\note{« Alg. »
est inscrit en interligne au-dessus du $\mathcal{B}$ de la question, que suit
une courbe}

\marginal{ok. dans le cas « classique » $X$ \struck{\ill{}} polycylindre
analytique \ldots}\note{au pied de la page, à droite, sous le crochet fermant
la question}

\page{8}

\note{feuillet d'un autre tapuscrit que celui des pages 2, 4 et 6, frappé sur
une machine sans accents, paginé « 8 » à la main en haut à droite : la fin
d'un paragraphe 6 sur les cycles premiers associés d'un Module quasi-cohérent.
La machine n'avait ni les flèches ni les signes ensemblistes ; il en a tracé
quelques-uns à la plume, et on rétablit les autres, avec les accents, là où le
texte les fixe}

\[
\mathrm{Ass}\, F = \bigcup_{i} \mathrm{Ass}\, F_{i}
\]
(iii)\quad $\mathrm{Ass} \coprod F_{i} = \bigcup_{i} \mathrm{Ass}\, F_{i}$ pour
toute famille $(F_{i})$ de modules qu\add{asi-}c\add{ohérents}.\note{les
réunions et la somme sont tracées à la main, la seconde réunion sur un « Ass »
tapé qu'elle surcharge ; le bout de ligne est coupé au bord du feuillet. Suit,
tapé : « Beweis klaar. »}

\textbf{Proposition 6.7.} Supposons $X$ loc noeth. Pour que la section $f \in
\Gamma(\mathcal{O}_{X})$ définisse une homothétie $F \to F$ qui est
\emph{injective}, il f et s que $X_{f} \supset \mathrm{Ass}\, F$, i.e. que $V(f)
\cap \mathrm{Ass}\, F = \emptyset$.\note{le $\Gamma$, la flèche, l'inclusion et
l'intersection manquent au tapuscrit et sont rétablis sur le sens}

\marginal{démontrer directement}\note{en regard de « En effet, on est ramené au
cas affine (par exemple) où c'est dans Bourbaki », avec un trait qui y renvoie}

\textbf{Corollaire 6.8.} Soit $I$ un Idéal cohérent sur $X$ loc noeth, $F$ un
Module cohérent sur $X$. Pour que $I$ soit « non diviseur de $0$ dans $F$ » il f
et s que $V(I) \cap \mathrm{Ass}\, F = \emptyset$.

\note{suit, tapé : « On peut mettre ici si on veut IV 2.8.18. et 2.8.19. (bien
qu'il soit possible qu'on ne les utilisera plus par la suite \ldots) »}

\textbf{Proposition 6.9.} Soient $X$ préschéma loc noeth, $U$ un ouvert dense.
Pour que $X$ soit réduit, il f et s que $U$ le soit, et que $X$ soit sans
« cycles premiers immergés »$^{*}$, ou encore que $X$ soit sans cycles premiers
imm. et que l'on ait $\mathrm{long}\, \mathcal{O}_{x} = 1$ pour $x$ point gén
d'une comp irr de $X$. (Cette notion $^{*}$ devrait être définie, pour un
Module, dans 6.1.)\note{le reste du feuillet, tapé, démontre la proposition et
observe : « Il est peut-être préférable d'essayer de donner de telles
démonstrations 'intuitives', plutôt que de se réduire toujours à l'énoncé connu
dans le cas affine », puis commence la définition d'un Module qu coh
\emph{réduit}}

\page{9}

\ill{} que le foncteur $M \mapsto \widetilde{M}$ des $\Gamma\mathcal{B}$-
\struck{$\mathcal{M}$} modules de présentation finie dans les
$\mathcal{B}$-Modules de prés. finie est pl. fidèle.
\struck{Prouvons \ill{}. Reste à prouver que \ill{}}

\add{$M$ \ill{} un} $\mathcal{B}$-Module de prés. finie, il est engendré par ses
sections. \struck{\ill{} $x$ \ill{} \ill{} \ill{}}\note{« $M$ \ldots\ un » est
porté au-dessus de la ligne biffée et relié par un trait vertical au début de
la ligne suivante}

\struck{N.B. $X$ \ill{} [(Nous supposons $X$ « pur » !)]. Alors
$\mathcal{M}_{x}$ est un $\mathcal{B}_{x}$-Module de présentation finie. Soit
$\mathcal{M}'$ le \ill{} \ill{} engendré par les sections.}\note{bloc encadré,
ouvert par « N.B. », et annulé par un long trait sinueux qui le traverse d'un
bout à l'autre et par deux obliques qui descendent de la ligne précédente}

(en utilisant seulement la condition $\alpha$) de la « remarque »
ci-dessus). Reste à prouver $\beta$) [qui est inclus dans le cas \ill{} \ill{}
\ill{} !], sur les conditions \ill{}. Il faut prouver que $\mathcal{B}$ satisfait
condition (ii) \& (iii), \uncertain{pour} \ill{} \ill{} (iii) \ill{} \ill{}
\ill{} \ill{} : la limite, \ill{}
\[
J = \varinjlim_{m} J_{m}
\]
il faut prouver que les $J_{m} = J \cap B_{m}$ \uncertain{sont} cohérents
\ldots

\marginal{? histoire d'anneaux cohérents ?}\note{dans le coin inférieur gauche,
écrit de biais}

\page{10}

\note{feuillet du tapuscrit des pages 2, 4 et 6, paginé « -- 101 -- » : la dimension des espaces topologiques. Les accolades
ensemblistes y sont tracées à la main, la machine ne les ayant pas}

\marginal{codim. \ill{} \ill{} \ldots}\note{en regard de la proposition 1,
écrit de biais, avec un trait qui y renvoie}

\textbf{Proposition 1.} (i) Si $Y$ est une partie d'un espace topologique $X$,
$\dim(Y) \leq \dim(X)$.

(ii) Si un espace topologique $X$ est réunion finie de parties fermées $X_{i}$,
alors $\dim(X) = \sup_{i}(\dim(X_{i}))$.\note{une troisième assertion (iii),
tapée, est biffée à la machine, et avec elle la ligne qui la démontrait. Suit
l'exemple d'une application continue $f$ avec $\dim(f(X)) > \dim(X)$ : $X$
discret à deux éléments $a, b$, $Y = \{a, b\}$ muni de la topologie dont les
fermés sont $\emptyset$, $\{a\}$ et $\{a, b\}$}

\textbf{Corollaire.} Soient $x \in X$, $U$ un voisinage de $x$, $Y_{i}$ ($1 \leq
i \leq n$) des parties fermées de $U$ telles que $x \in Y_{i}$ pour $1 \leq i
\leq n$ et que $U$ soit la réunion des $Y_{i}$. Alors $\dim_{x} X =
\sup_{i}(\dim_{x}(Y_{i}))$.

\textbf{Proposition 2.} (i) Soit $\Phi$ l'ensemble des parties fermées
irréductibles de $X$. On a $\dim(X) = \sup_{Y \in \Phi}(\mathrm{codim}_{X}(Y))$.

(ii) On a $\dim(X) = \sup_{x \in X} \dim_{x}(X)$.

\page{11}

\textbf{\struck{Corollaire} \struck{Proposition} Corollaire}\note{le titre est
écrit trois fois, souligné chaque fois : « Corollaire », biffé et surchargé de
« Proposition », elle-même biffée, puis « Corollaire » au-dessus} Supposons que
$X$ satisfasse la condition \struck{\ill{}} suivante plus forte que (ii) :
\begin{itemize}
\item[(ii')] (\add{les $\underline{O}_{X,x}$ noeth., et}) Pour tout entier $n$,
$\underline{O}_{X}[T_{1}, \ldots, T_{n}]$ est un faisceau d'anneaux
cohérent.\note{l'ajout est entouré et renvoyé par un trait au début de la
condition ; un grand point d'interrogation est porté en marge en regard}
\item[(iii')] Tout $x \in X$ a un syst. fond. de voisinages quasi-compacts
satisfaisant la condition (iii) \uncertain{pour tout} faisceau \ill{} sur
$\underline{O}_{X}$ \ldots\note{la condition est ajoutée en petits caractères
entre les lignes, sa fin sur deux niveaux}
\end{itemize}
Alors la conclusion du \uncertain{Lemme} A reste valable en y remplaçant
$\underline{O}_{X}$ par $\underline{O}_{X}[T_{1}, \ldots, T_{n}]$, plus
généralement par n'importe quel faisceau d'Algèbres $\mathcal{B}$ de
présentation finie.

Soit
\[
\mathcal{B}^{p} \xrightarrow{u} \mathcal{B}^{q} \to \mathcal{M} \to 0
\qquad (u = (u_{ij}),\ (i, j) \in P \times Q)
\]
une suite exacte, $\Sigma$ \struck{\ill{} \ill{}} de $\mathcal{B}$-Modules
($\mathcal{B} = \underline{O}_{X}[T_{1}, \ldots, T_{n}]$) \ill{} \ill{} pour tt
\ill{} une \uncertain{sous-suite} exacte $\Sigma_{m}$ de $\underline{O}_{X}$-Modules
cohérents
\[
\mathcal{B}_{m}^{p} \to \mathcal{B}_{N+m}^{q} \to \mathcal{M}_{m} \to 0
\]
($N$ majorant des degrés des $u_{ij}$) \struck{\ill{}} et $\Sigma$ est limite
des $\Sigma_{m}$. \uncertain{On a donc}
\[
\Gamma\mathcal{B}_{m}^{p} \to \Gamma\mathcal{B}_{m+N}^{q} \to \Gamma\mathcal{M}_{m} \to 0
\]
d'où en passant à la limite
\[
\Gamma\mathcal{B}^{p} \to \Gamma\mathcal{B}^{q} \to \Gamma\mathcal{M} \to 0
\]
ce qui implique $\alpha$. \struck{\ill{}} Ceci prouve \ldots\note{le
$\alpha$ renvoie à la condition ($\alpha$) de la remarque de la page 5 ; la
page s'arrête sur « Ceci prouve », et le haut de la page 9, « \ldots\ que le
foncteur $M \mapsto \widetilde{M}$ \ldots\ est pl. fidèle \ldots\ (en
utilisant seulement la condition $\alpha$) », semble en être la suite ; si
c'est le cas, l'ordre des archivistes n'est pas ici celui de la rédaction}

\page{12}

\note{feuillet du même tapuscrit, paginé « -- 102 -- », qui
fait suite à la page 10 : fin de la démonstration de la proposition 2,
corollaires 1 à 4, proposition 3. Il n'y porte à la main que des signes
ensemblistes et un renvoi, et on ne donne que les énoncés}

\textbf{Corollaire 2.} Soit $X$ un espace noethérien vérifiant l'axiome
$(T_{0})$ et soit $F$ l'ensemble des points fermés de $X$. Alors $\dim(X) =
\sup_{x \in F}(\dim_{x}(X))$.

\textbf{Corollaire 3.} Soit $Y$ une partie fermée de $X$. Alors
\[
(1) \qquad \dim(Y) + \mathrm{codim}_{X}(Y) \leq \dim(X).
\]

\textbf{Corollaire 4.} Soient $Y, Z, T$ trois parties fermées irréductibles de
$X$ telles que $Y \subset Z \subset T$. Alors
\[
(2) \qquad \mathrm{codim}_{T}(Y) \geq \mathrm{codim}_{Z}(Y) +
\mathrm{codim}_{T}(Z).
\]

\textbf{Proposition 3.} Soit $X$ un espace noethérien vérifiant l'axiome
$(T_{0})$. Pour que $\dim(X) = 0$, il faut et il suffit que $X$ soit fini et
discret. Pour qu'une partie fermée $Y$ de $X$ soit de codimension nulle dans
$X$, il faut et il suffit que $Y$ contienne une composante irréductible de
\add{$X$}.\note{à la dernière ligne de la démonstration, il insère au-dessus de
« toute composante irréductible de $X$ contient un point fermé » le renvoi
« (§ 4, n° 1, lemme 1) »}

\page{13}

\textbf{\struck{Prop.} Lemme B}\quad Supposons \struck{que} que $X$ satisfait
les conditions suivantes : (i) (ii) et
\begin{itemize}
\item[(iv)] Les $\underline{O}_{X,x}$ ($x \in X$) sont \emph{noethériens}
\item[(v)] Pour tout \struck{\ill{}} Idéal coh. $J$, \struck{\ill{}} posant
$Y = (\mathrm{Supp}\, \underline{O}_{X}/J, \underline{O}_{X}/J)$, et tout
$y \in Y$ tel que $\underline{O}_{Y,y}$ est \emph{intègre}, il existe un
voisinage ouvert $U$ de $y$ \add{dans $Y$} ayant la propriété suivante : toute
section $f$ de $\underline{O}_{Y}$ sur \struck{un voisinage} $U$ qui est
\struck{\ill{} \ill{}} nulle au voisinage de $y$, est nulle sur
$U$.\note{« dans $Y$ » est en interligne, relié par un trait ; la lettre qui
précède le signe « $=$ » dans la définition de $Y$ est surchargée}
\end{itemize}
Alors pour tout \struck{\ill{}} $F$ coh. sur $X$, et toute suite croissante
$(F_{i})$ de sous-faisceaux coh. de $F$ est stationnaire \ldots\note{« croissante »
est inscrit en interligne}

\textbf{Corollaire}\quad Sous les conditions (i) à (v), $A = \Gamma(X,
\underline{O}_{X})$ est un anneau \emph{noethérien}.

Ces conditions sont satisfaites si $X$ est un \struck{espace de Stein}
\add{polycylindre analytique} (compact), \struck{\ill{}} sur $\mathbb{C}$
\ill{} \ill{}. \struck{« \ill{} \ill{} \ill{} »} et les communs
\ldots\note{« polycylindre analytique » est écrit au-dessus de « espace de
Stein », biffé ; les mots qui suivent « (compact) » sont surchargés et en
partie encadrés, et la dernière ligne, entre guillemets, est barrée d'un trait}

\marginal{premier espoir}\note{entouré, dans la marge inférieure gauche, en
regard du dernier alinéa}

\page{14}

\note{feuillet du même tapuscrit, paginé « -- 104 -- » :
la proposition 5 sur la condition des chaînes, et le début de sa démonstration.
Il y est intervenu plus que sur les autres, en substituant partout la notion
d'espace équidimensionnel à celle de points fermés de même codimension ; les
passages biffés à la machine ne sont pas relevés}

\textbf{Proposition 5.} Soit $X$ un espace noethérien, de dimension finie,
vérifiant l'axiome $(T_{0})$. Les conditions suivantes sont équivalentes :

a) Deux chaînes maximales de parties fermées irréductibles de $X$ ont même
longueur.

b) \add{$X$ est équidimensionnel et} \struck{\ill{}, les points fermés de $X$
ont même codimension}, et $X$ satisfait à la condition des chaînes.\note{« $X$
est équidimensionnel et » est écrit au-dessus de la ligne tapée ; « et
équidimensionnel » est d'autre part porté en marge en regard, avec un renvoi}

c) $X$ est équidimensionnel et pour deux parties fermées irréductibles $Y
\subset Z$ de $X$, on a
\[
(4) \qquad \dim(Z) = \dim(Y) + \mathrm{codim}_{Z}(Y).
\]

d) \add{$X$ est équidimensionnel} \struck{Les codimensions des points fermés de
$X$ sont les mêmes}, et pour deux parties fermées irréductibles $Y \subset Z$
de $X$, on a
\[
(5) \qquad \mathrm{codim}_{X}(Y) = \mathrm{codim}_{Z}(Y) +
\mathrm{codim}_{X}(Z).
\]
\note{la substitution est en marge, avec un renvoi à « d) »}

\marginal{\emph{Dim} \quad \emph{Cor.} Soit $X$ un esp. \ill{} \ill{} irréd. de
$X$, et $x$ un pt fermé de $X$. \ill{} $\dim X = \mathrm{codim}_{X} x =
\uncertain{\dim_{X} x}$ \ill{}}\note{écrit de biais dans la marge gauche, en regard de la
démonstration ; les deux derniers membres de l'égalité sont lus sur les
lettres et ne s'accordent pas avec ce qu'on attendrait}

Dans la démonstration, il précise « toute chaîne saturée
\add{d'extrémités} $Y$ et $Z$ \add{$(Y \subset Z)$} est contenue dans une
chaîne maximale dont les éléments \add{autres que ceux de la chaîne donnée}
sont ou bien contenus dans $Y$ ou bien contiennent $Z$ », puis « prouvent aussi
que \add{si a) est vérifiée}, pour toute partie fermée irréductible $Y$ de
$X$, on a
\[
(6) \qquad \dim(Y) + \mathrm{codim}_{X}(Y) = \dim(X) \text{ »,}
\]
formule qu'il entoure, et corrige ensuite les numéros des formules qui s'en
déduisent et ajoute le renvoi « (prop. 4) ».\note{les deux numéros corrigés
sont surchargés et ne se lisent pas avec certitude}

\section*{Un tore complexe qui n'est pas une variété abélienne}

\note{titre de l'édition : les deux feuillets qui suivent, écrits à l'encre
noire d'une main posée, ne portent pas de titre, et rien ne les rattache aux
pages 1 à 14}

\page{16}

Il existe un tore complexe, de dim (complexe) 2, fibré sur un tore complexe de
dimension 1, \struck{à fibres} par un sous-tore complexe de dim 1, qui n'est pas
une variété abélienne. (À fortiori, un fibré analytique à base et fibres des
variétés de Hodge, n'est pas nécessairement une variété de Hodge)

\emph{Exemple} : On prend le tore $T$ quotient de $\mathbb{C}^{2}$ par le groupe
$\Gamma$ engendré par
\[
e_{1}, \quad ie_{1}, \quad e_{2}, \quad e_{1} + ike_{2}
\]
avec $k$ \emph{irrationnel}, et le sous-tore $T_{0}$ image de l'axe complexe
$V$ engendré par $e_{1}$ (donc $T_{0} = \mathbb{C}/\mathbb{Z}^{2}$). Si $T$
était abélienne, on sait qu'il existerait un supplémentaire (complexe) $V'$ de
$V$, tel que $V$ soit $\Gamma$-\uncertain{propre}, i.e. $V$ engendré par $\Gamma
\cap V$.\note{c'est bien $V$, et non $V'$, qu'il écrit aux deux endroits ; la
suite du raisonnement porte sur $V'$} À fortiori, il existerait $x \in
\struck{\Gamma \cap V} V'$ \struck{tel que $ix$ soit de la forme $\rho x +
\sigma y$, $\rho$ et $\sigma$ réels, $y \in \Gamma \cap V$} \ill{} $u$ complexe,
$ux \in \Gamma$, (prendre $x \in \Gamma \cap V$ non nul). Le
« \uncertain{minimum} »\note{le bas de la page est un bloc surchargé : deux
lignes barrées d'un trait, une correction « $\rho x + \sigma ux$ » portée
au-dessus, le tout traversé d'un trait sinueux et fermé à gauche par un trait
vertical. On donne ce qui se lit ; la reprise est au haut de la page suivante}

\page{17}

et un nb complexe \emph{non réel} $c = c' + ic''$ ($c', c'' \in \mathbb{R}$) tel
que $x$ et $cx \in \Gamma$. (resp. $x \in V' \cap \Gamma$, $y \in V' \cap
\Gamma$ \struck{linéairement} lin. indép. sur $\mathbb{R}$) Les composantes de
$x$ seront
\[
\alpha + \delta, \quad \beta, \quad \gamma, \quad k\delta, \qquad (\gamma
\text{ ou } \delta \neq 0 ;\ \alpha\beta\gamma\delta \in \mathbb{Z})
\]
celles \struck{$\beta\gamma\delta$} de $cx$ seront
\[
c'(\alpha + \delta) - c''\beta, \quad c''(\alpha + \delta) + c'\beta, \quad
c'\gamma - k\delta c'', \quad c''\gamma + c'k\delta
\]
Il faudrait donc que ces dernières valeurs puissent s'écrire
$\alpha' + \delta', \beta', \gamma', k\delta'$, avec $\alpha'\beta'\gamma'\delta'
\in \mathbb{Z}$. En particulier, les premières 2 composantes devraient être
entières, d'où on conclut que $c'$ et $c''$ seraient \emph{rationnels}. Pour que
les deux \struck{dernières} \uncertain{troisièmes} composantes puissent être
rationnelles, il faut donc que $k\delta c'' = 0$ \struck{\ill{}}. Comme $c''
\neq 0$, cela implique $\delta = 0$. Donc les composantes de $cx$ sont
\[
\alpha c' - \beta c'', \quad \beta c' + \alpha c'', \quad c'\gamma, \quad
c''\gamma
\]
\ill{} \ill{} $k\delta' = c''\gamma$, ce qui n'est possible que si les deux
valeurs sont nulles, d'où $\gamma = 0$. Mais, \struck{donc} \add{comme $\gamma =
\delta = 0$}, $x$ serait dans l'espace engendré par $e_{1}$ et $ie_{1}$,
absurde.\note{« comme $\gamma = \delta = 0$ » est inscrit au-dessus de « donc »
biffé, et « engendré par » au-dessus de la dernière ligne}

\end{document}
